\subsection{近代物理实验}

\begin{结论}{电子自旋共振测 $g$ 因子}
共振条件 \重点{h\nu=g\mu_BB}，$\frac{\mu_B}{h}=13.996\,$GHz/T。固定微波（或射频）频率扫磁场，测共振场 $B_0$：$g=\frac{h\nu}{\mu_BB_0}$。
标样 DPPH $g=2.0036$；自由电子 $g_e=2.0023$。场调制 $+$ 锁相检测得到吸收线的一阶导数，峰峰间距对应线宽（$\Delta B_{pp}=\frac{2}{\sqrt3}\cdot$ 半高半宽，Lorentz 线型）。
\end{结论}

\begin{结论}{Mössbauer 谱：用 Doppler 速度扫能量}
无反冲共振吸收，源以速度 $v$ 运动：\重点{\frac{\Delta E}{E_\gamma}=\frac vc}。${}^{57}$Fe：$E_\gamma=14.4\,$keV，自然线宽 $\Gamma=4.7\times10^{-9}\,$eV，对应 $v=c\Gamma/E_\gamma\approx0.1\,$mm/s。
$\alpha$-Fe 磁超精细分裂给六条线（外侧间距约 $10.6\,$mm/s）；同质异能移位反映核处电子密度。Pound--Rebka 用它测出引力红移 $gh/c^2\approx2.5\times10^{-15}$（$h=22.5\,$m）。
\end{结论}

\begin{结论}{$\gamma$ 能谱的结构}
单能 $E$ 的 $\gamma$ 射线在闪烁体/半导体中：全能峰；Compton 坪及其边缘 $E_C=\frac{2E^2}{m_ec^2+2E}$；背散射峰 $E-E_C=\frac{E}{1+2E/m_ec^2}$。
${}^{137}$Cs（$662\,$keV）：边缘 $478\,$keV，背散射峰约 $184\,$keV。$E>1.022\,$MeV 时出现单、双逃逸峰（$E-511$、$E-1022\,$keV）。分辨率：NaI 约 $7\%$，HPGe 约 $0.2\%$（$662\,$keV）。
\end{结论}

\begin{结论}{$\mu$ 子寿命实验}
停止在闪烁体中的宇宙线 $\mu$ 子，测“停止信号”与“衰变电子信号”的时间差直方图，拟合 $N(t)=N_0\ee^{-t/\tau}+B$（$B$ 为偶然符合本底）。
\[
\重点{\tau_{\mu^+}=2.197\,\mu\mathrm s},\qquad \mu^-\text{ 在碳中因核俘获缩短为约 }2.03\,\mu\mathrm s .
\]
混合样品得到两指数之和，单指数拟合约 $2.1\,\mu$s。由 $\tau_\mu$ 与 $m_\mu$ 可推出 Fermi 常数 $G_F$。
\end{结论}

\begin{结论}{Rb 光学抽运：用 Zeeman 共振测弱磁场}
圆偏振 $D_1$ 光把原子抽运到边缘 $m_F$ 态，样品变透明；射频频率满足 Zeeman 共振时布居被打乱、透射下降：
\[
\重点{\nu=g_F\frac{\mu_BB}{h}},\qquad {}^{87}\mathrm{Rb}\ (g_F=\tfrac12):\ 0.700\,\mathrm{MHz/G};\quad {}^{85}\mathrm{Rb}\ (g_F=\tfrac13):\ 0.467\,\mathrm{MHz/G}.
\]
两同位素共振频率之比 $3/2$ 可确认谱线归属。扫亥姆霍兹线圈电流，零场截距给出地磁场分量（约 $0.5\,$G，即 $350\,$kHz）。
\end{结论}

\begin{结论}{Compton 实验的线性化}
不同散射角 $\theta$ 测散射光子能量 $E'$：
\[
\重点{\frac1{E'}-\frac1E=\frac{1-\cos\theta}{m_ec^2}},
\]
作 $1/E'$ 对 $1-\cos\theta$ 的直线，斜率给 $1/m_ec^2$（应得 $511\,$keV），截距应为 $1/E$。计数率随角度的变化与 Klein--Nishina 截面比较。
\end{结论}

\begin{结论}{Balmer 系与氘的同位素位移}
\[
\frac1\lambda=R_M\Bigl(\frac14-\frac1{n^2}\Bigr),\qquad \重点{R_M=\frac{R_\infty}{1+m_e/M}}.
\]
H$_\alpha$ $656.28\,$nm，D$_\alpha$ $656.10\,$nm，差 $0.18\,$nm——Urey 由此发现氘。用光栅测 $\lambda_{3\to2},\lambda_{4\to2},\dots$，作 $1/\lambda$ 对 $1/n^2$ 的直线，斜率 $-R_M$。
\end{结论}

\begin{结论}{钨灯验证 Stefan--Boltzmann}
钨丝温度由电阻比确定（$\rho\propto T^{1.2}$ 附近，查表更准），电功率 $P=VI$ 近似等于辐射功率（高温时导热损失可忽略）：
\[
\ln P=4\ln T+\text{常数}\quad\Rightarrow\quad \重点{\text{双对数图斜率}\approx4}.
\]
低温段斜率偏离 4 说明传导/对流损失或发射率随温度变化；钨发射率约 $0.3$--$0.45$，非理想黑体。
\end{结论}

\begin{结论}{Rutherford 散射计数率的角度依赖}
探测器固定立体角，计数率 $\propto\frac{1}{\sin^4(\theta/2)}$。作 $\ln N$ 对 $\ln\sin\frac\theta2$ 的直线，斜率应为 \重点{-4}。
小角度处偏离来自束流发散与探测器有限角宽（需卷积）；大角度计数极少，需长时间累积并扣本底。改变箔厚检验 $N\propto t$（单次散射近似）；改变靶材检验 $N\propto Z^2$。
\end{结论}

\begin{结论}{Geiger 计数的 Poisson 检验}
固定时间窗重复计数 $k$ 次：Poisson 要求 $\bar n\approx s^2$。离散指数
\[
\重点{D=\frac{(k-1)s^2}{\bar n}\sim\chi^2_{k-1}}.
\]
$D$ 显著偏大：源不稳、计数窗漂移；偏小：死时间（计数被“规则化”）。高计数率下先做死时间修正再检验。分布形状用 $\chi^2$ 拟合优度（合并期望数 $<5$ 的格）。
\end{结论}

\begin{结论}{饱和吸收光谱：Doppler-free 分辨}
强泵浦光与弱探测光反向穿过气室：只有 $v_z=0$ 的原子同时与两束共振，探测透射在跃迁频率处出现窄峰（Lamb 凹陷），宽度接近自然线宽。
两跃迁 $\nu_1,\nu_2$ 共享基态时，在 \重点{\frac{\nu_1+\nu_2}{2}} 处出现交叉共振峰（速度 $v_z=\pm\frac{c(\nu_2-\nu_1)}{2\nu}$ 的原子）。用于激光稳频与 Rb 超精细结构测量。
\end{结论}

\参考{Melissinos--Napolitano, \emph{Experiments in Modern Physics}；MIT 8.13--8.14 Junior Lab 实验讲义（ESR、Mössbauer、$\gamma$ 谱、$\mu$ 子寿命、光抽运、Compton、Rutherford）；Knoll, \emph{Radiation Detection and Measurement}；Demtröder, \emph{Laser Spectroscopy}。}
